\label{sec:conclusao}

Este projeto propôs medir em que proporção os SDKs de rastreamento de terceiros embarcados nos aplicativos Android governamentais brasileiros estão declarados ao cidadão, tomando como unidade o par aplicativo--SDK e como referência a LGPD. Com base em estudos que encontraram declarações ausentes ou incorretas em parte relevante dos aplicativos analisados, a hipótese é que ao menos 10\% desses pares não terão declaração compatível na política de privacidade nem na seção ``Segurança dos dados'' da Google Play.

No Projeto Capstone I, foram definidos o problema, os objetivos, a hipótese e o método, e foram realizadas as etapas preparatórias: a identificação dos 376 aplicativos classificados como oficiais, a consulta ao AndroZoo, o piloto com 20 pacotes e a revisão sistemática da literatura, cuja leitura completa está em andamento. Esses resultados indicam que o método é viável e já orientaram decisões importantes, como a coleta na Google Play e a comparação entre detectores, mas não antecipam a decisão sobre a hipótese.

As principais limitações são o tipo de evidência, já que um SDK presente no código não prova envio de dados; o caráter das declarações, que tratam de categorias de dados e finalidades, e não da lista de SDKs; e a restrição ao Android. Por isso, o tráfego será observado em uma subamostra, a ausência do nome de um SDK não será tratada como omissão, e as versões iOS ficam como trabalho futuro.

No Projeto Capstone II, o cronograma da Tabela~\ref{tab:cronograma} prevê concluir a revisão sistemática, revisar a lista de aplicativos, baixar e analisar os pacotes, classificar as declarações, observar o tráfego da subamostra e decidir sobre a hipótese. Espera-se, com isso, oferecer uma medida reprodutível do quanto os aplicativos públicos declaram o que neles está embarcado e subsídios para que os órgãos responsáveis revisem suas declarações à luz da LGPD.
